% Tuning: one machine parameter -> one code choice, and the result of all of them together.

The analysis tied each effect to one parameter of this machine. This section turns each
parameter into a rule that \code{scan\_bench} applies when the corresponding knob is set, and
measures what it buys (\cref{tab:tuning}). The rules are computed at start-up from sysfs, so
nothing in the code is specific to this laptop; the constants of the materialization threshold
(\cref{sec:mat}) were fitted on it.

\begin{table*}[t]
\centering\small
\caption{Machine parameters and the choices derived from them.}
\label{tab:tuning}
\begin{tabularx}{\textwidth}{>{\raggedright}p{3.3cm}>{\raggedright}p{2.5cm}Xp{3.7cm}}
\toprule
Parameter (this machine) & Decides & Rule & Measured effect \\
\midrule
Line size (64\,B), value width; where the table lives & early or late, per span &
late while $f(s) < \varphi(1-\sigma/g)$, with constants for DRAM and for tables that fit in L3
(\cref{sec:mat:auto}) & within 1\% of the better fixed strategy on DSM from DRAM; \code{SELECT *}
in L3 worst case 2.6$\times$ $\to$ 1.3$\times$ \\
Cache sizes vs work per value & vector size & fit the vector's buffers in L2 ($\frac14$) for
our scans; smaller for random access inside a vector, larger for a high fixed cost per vector &
decoding SUM 1.2$\times$ faster than an L1-sized vector (\cref{sec:tuning:batch}) \\
L1 set stride (4\,KB), 64 sets, 8 ways & DSM column starts, PAX minipage stride & column $c$
starts $17c$ lines into its page; minipages hold an odd number of lines & early \code{SELECT *}
at 1\%: 2.7$\times$ on DSM, 2.3--2.7$\times$ on PAX (\cref{sec:tuning:place}) \\
Write-allocate (12\,GB/s, 40 with NT stores) & how output is written & build each vector's
tuples in a cache-resident stage, copy it out with streaming stores & projections of every row
1.2--1.9$\times$ faster (\cref{sec:tuning:output}) \\
Prefetcher (stops at 4\,KB pages); dTLB reach (8\,MB) & PAX page size; OS page size & 256\,KB
PAX pages; 2\,MB pages for buffers of 2\,MB or more & PAX SUM(c0) 4.5 $\to$ 1.2\ms; late SUM
20\% faster (\cref{sec:tuning:pages}) \\
AVX2 (no compress, no AVX-512); DRAM at 7\,B/cycle & selection and decoding kernels; code width
& hand-written SIMD selection vector and fused unpack-and-add; dictionary codes of 1, 2 or 4
bytes & late SUM wins in cache; bit-packed SUM 1.6$\times$ faster than plain; dictionary
predicate 2.2--5.4$\times$ (\cref{sec:mat:cache,sec:enc}) \\
\bottomrule
\end{tabularx}
\end{table*}

\subsection{Vector size: match it to the work, not to one cache level}
\label{sec:tuning:batch}

\begin{table}[t]
\centering\small
\caption{ns per value for three synthetic workloads and our decoding scan, with the vector's
buffers sized to fit L1, L2 or L3 (\code{build/batchbench}, \code{results/batch\_demo/}). Best
per row in bold.}
\label{tab:batchwork}
\begin{tabular}{lrrr}
\toprule
Workload & fits L1 & fits L2 & fits L3 \\
\midrule
random reads inside the vector & \textbf{1.34} & 1.69 & 2.55 \\
random reads + 2\,\textmu s per vector & 2.40 & \textbf{1.90} & 2.60 \\
light work + 20\,\textmu s per vector & 12.05 & 2.26 & \textbf{1.69} \\
our scan (decode 8 bit-packed columns + SUM) & 0.37 & \textbf{0.31} & 0.32 \\
\bottomrule
\end{tabular}
\end{table}

The engine processes one vector of $B$ rows at a time, and the vector's buffers (selection
vector, decoded values, output tile) live in L1, L2 or L3 depending on $B$. Two costs pull in
opposite directions: the work per value, which is cheaper when the buffers sit in a faster
level, and the fixed cost per vector (a decode call per column, loop setup, a prefetcher
restart, or in other systems a call into another runtime or a hand-off to another thread),
which a larger vector amortises. \Cref{tab:batchwork} shows that either can win.
\begin{itemize}
\item \textbf{Random reads inside the vector} (like probing a per-vector table) want L1: L1
answers in about 4 cycles, L2 in 12.
\item \textbf{The same with 2\,\textmu s per vector} wants L2: an L1-sized vector pays the fixed
cost too often.
\item \textbf{Light work with 20\,\textmu s per vector} wants L3 or more: the fixed cost
dominates until the vector is megabytes large.
\item \textbf{Our scans} read their buffers sequentially, and an L2 hit streams at 130\,GB/s,
nearly L1's 122, so L1-sized vectors cut L1D misses but not time. Late \code{SELECT *} at
100\% with 128-row vectors keeps its output tile in L1 and cuts misses from 19 to 5 per row,
at the same latency. But small vectors slow decoding (one call per column per vector):
bit-packed SUM takes 11.9\ms{} with the 128 rows an L1-sized rule gives, and 9.9\ms{} with
1024. Beyond the L2 the buffers spill and both slow down again.
\end{itemize}
So \code{--batch auto} sizes the vector so that every byte it touches per row fits in a quarter
of L2, $B = \max\{2^i : 2^i \cdot b_\text{row} \le L2/4\}$, leaving room for lines in flight and
prefetched lines. On this machine that is 256 rows for \code{SELECT *} of 16 attributes, 1024
for 4, and 1024 for decoding 8 bit-packed columns. A workload with random accesses inside the
vector, or a large fixed cost per vector, would need the rule moved to L1 or to L3.

\subsection{Placing columns and minipages in the cache sets}
\label{sec:tuning:place}

\begin{table}[t]
\centering\small
\caption{Early \code{SELECT *} at 1\% on PAX, plain minipages vs an odd number of lines per
minipage (\code{--pax-pad auto}); L1D misses per row in brackets.}
\label{tab:aliasing}
\begin{tabular}{lrr}
\toprule
Page size (minipage stride) & plain & odd-line minipages \\
\midrule
32\,KB (2\,KB: 8 per set, fits) & 37.8\ms & not needed \\
64\,KB (4\,KB) & 65.6\ms{} (8.6) & 27.8\ms{} (3.3) \\
256\,KB & 63.0\ms & 24.2\ms \\
4\,MB & 62.7\ms & 23.3\ms \\
\bottomrule
\end{tabular}
\end{table}

Streams read in lockstep must not share an L1 set (\cref{sec:analysis:cache}). Two rules
follow from the 64 sets and the 4\,KB set stride.
\begin{itemize}
\item \textbf{DSM} (\code{--dsm-offset auto}): column $c$ starts $17c$ lines into its
page-aligned allocation. 17 is coprime to the 64 L1 sets and the 1024 L2 sets, so up to 64
columns fall into distinct L1 sets. Early \code{SELECT *} at 1\% drops from 62.6 to 22.9\ms.
\item \textbf{PAX} (\code{--pax-pad auto}): from 64\,KB pages up, the 16 minipages start $P/16$
apart, a multiple of 4\,KB, so value $r$ of all of them maps to one set. Shortening each
minipage to an odd number of lines (64 $\to$ 63 at 64\,KB) puts them in 16 different sets, at a
cost of one line per minipage (1.6\%). Early \code{SELECT *} at 1\% becomes 2.3--2.7$\times$
faster, within 7\% of DSM from 256\,KB pages (\cref{tab:aliasing}). Padding by 64 bytes does
not help: with the page size fixed, the stride stays $P/16$.
\end{itemize}
The effect is largest where reading dominates. Early is the only strategy that reads the 16
streams at once (late reads one at a time and is unaffected), and at 1\% almost nothing is
written. At 100\%, where writing the output dominates, the misses still halve but the time
improves only 1.1$\times$ (90.8 to 81.9\ms{} at 64\,KB), and SUM of 16 attributes, which is
bound by DRAM bandwidth, also gains only 1.1$\times$ (20.0 to 18.4\ms).

We found the effect by accident. DSM columns first came from \code{malloc}, and depending on
the allocator's state they were either carved back to back out of the heap (which happened to
spread them over 16 sets) or given one page-aligned \code{mmap} each. The same \code{SELECT *}
took 21\ms{} in one experiment and 64\ms{} in another. The engine now gives large buffers
their own \code{mmap} and staggers DSM columns explicitly.

\subsection{Writing the output}
\label{sec:tuning:output}

\begin{table}[t]
\centering\small
\caption{Three ways to write a projection's output.}
\label{tab:stage}
\begin{tabular}{lcc}
\toprule
Approach & absorbs scattered writes? & avoids the read before write? \\
\midrule
normal stores & yes (the cache does) & no, paid per line \\
streaming stores directly & no (partial lines) & yes \\
stage (normal stores) + streaming copy & yes & yes \\
\bottomrule
\end{tabular}
\end{table}

\begin{table}[t]
\centering\small
\caption{Projections of every row (ms) with direct writes, a stage copied with \code{memcpy}, and
a stage copied with streaming (NT) stores.}
\label{tab:outputknobs}
\resizebox{\linewidth}{!}{\gen{output_knobs}}
\end{table}

A projection of every row writes as many bytes as it reads, and from DRAM those writes cost
more: a normal store that misses reads the line first, so writes reach 12\,GB/s while reads
reach 28. Streaming (non-temporal) stores skip that read: they fill a 64-byte write-combining
buffer in the core and send the whole line to memory, at 40\,GB/s. But only whole lines written
in order benefit, and our writes are scattered: early writes each row's slot and overwrites it
when the row fails, late writes a vector's output one attribute at a time. So
\code{--nt-stores} lets each vector write its tuples with normal stores into a small stage that
stays in cache, and copies the stage out with 32-byte streaming stores (\cref{tab:stage}).
\begin{itemize}
\item At 100\% this makes every configuration 1.2--1.9$\times$ faster
(\cref{tab:outputknobs}): DSM early 71.2 to 37.9\ms, NSM early 67.1 to 41.5\ms, and PAX late,
whose column-at-a-time writes were the most scattered, 86.5 to 45.3\ms. The line model predicts
it: 512\,MB read at 28\,GB/s and written at 40\,GB/s take 31\ms{} if nothing overlaps, against
61\ms{} with write-allocate.
\item The stage alone (\code{memcpy}) is up to 36\% slower than writing directly: it adds a
copy and still pays for write-allocate. The gain comes from the streaming stores.
\item At 1\% there is almost nothing to write, and the three modes are within noise.
\end{itemize}
Streaming stores bypass the cache, so they suit a result written once and read later. A result
consumed right away by the next operator should stay in cache.

\subsection{Pages: the prefetcher and the TLB}
\label{sec:tuning:pages}

\paragraph{PAX page size.} Every PAX minipage is a new stream that the prefetcher has to find
(\cref{sec:analysis:cache}), about 120\,ns per page at 16\,KB. Larger pages make the minipages
longer: SUM(c0) takes 4.9\ms{} at 16\,KB, 2.8 at 64\,KB and 1.2 at 4\,MB, the same as DSM. The
tuned configuration uses 256\,KB pages with odd-line minipages, where late \code{SELECT *}
reaches DSM late and early comes within 7\% of DSM early.

\paragraph{Huge pages.} The 2048-entry L2 dTLB covers 8\,MB with 4\,KB pages, so a 32\,MB column
needs 8192 translations, and 16 with 2\,MB pages. With \code{--hugepages}, buffers of 2\,MB or
more are allocated 2\,MB-aligned and marked \code{MADV\_HUGEPAGE} before they are first touched
(WSL2's kernel backs only such marked regions with huge pages). For SUM of 16 DSM columns, dTLB
misses fall from 0.031 to 0.008 per row for both strategies, but only late gets faster: 23.3 to
18.7\ms{} (20\%), against 19.3 to 18.8\ms{} (3\%) for early. Early's 16 interleaved columns keep
other loads going while one of them waits for a page walk; late works through one column at a
time, so its single stream stalls on every walk. Staggering helps early, huge pages help late.

\subsection{Everything together}
\label{sec:tuning:all}

\begin{figure}[t]
\centering
\fig{tuned}
\caption{Six queries with the untuned defaults (light) and with every machine-derived knob
(dark), best strategy of each, 4M rows from DRAM.}
\label{fig:tuned}
\end{figure}

\begin{table*}[t]
\centering\small
\caption{Untuned against tuned, in ms. Untuned: 1024-row vectors, page-aligned DSM columns,
4\,KB pages, direct output writes, scalar kernels, 16\,KB PAX pages, and the best of the four
fixed strategies. Tuned: \code{--batch auto --dsm-offset auto --hugepages --nt-stores --simd},
256\,KB PAX pages with \code{--pax-pad auto}, and auto (in this run measuring both strategies,
\cref{sec:tried}). Speedups are over the untuned best.}
\label{tab:tuned}
\gen{tuned}
\end{table*}

\Cref{fig:tuned,tab:tuned} put the knobs together on six queries and three layouts, from DRAM.
The untuned engine is given the best of the four fixed strategies per cell, which a real engine
cannot know in advance; the tuned one has to pick its strategy with auto.
\begin{itemize}
\item \textbf{Over the 18 cells the tuned engine is 1.46$\times$ faster} (geometric mean;
1.50$\times$ with the best tuned strategy): 1.11$\times$ on NSM, 1.41$\times$ on DSM and
1.96$\times$ on PAX, up to 3.8$\times$ for PAX SUM(c0) (4.5 to 1.2\ms).
\item \textbf{Projections of many rows gain on every layout}: \code{SELECT *} of every row
1.7$\times$ on NSM, 2.3$\times$ on DSM and 2.5$\times$ on PAX, mostly from the streaming stores.
After tuning NSM no longer wins it: 36.9\ms, against 34.6 on DSM and 31.7 on PAX.
\item \textbf{PAX gains the most}, because two of its defaults are wrong for memory: 16\,KB pages
are too short for the prefetcher, and larger pages alias in L1 without the odd-line stride.
\item \textbf{Queries that already move few lines gain little}: DSM SUM(c0) is one sequential
stream (1.03$\times$).
\item \textbf{Exceptions.} NSM \code{SELECT *} at 10\% (0.95$\times$): its best untuned
strategy, early-pipe, is not in the tuned set. PAX at 10\% is unchanged. One cell, NSM
\code{SELECT c0..c3} at 1\% with auto (0.76$\times$), hit a slow period on the host; repeated
three times, auto takes 15.2--16.1\ms{} against late's 14.5--15.4, and the 1.46$\times$ includes
the bad measurement.
\end{itemize}
